\usepackage{amsmath,amssymb}
\usepackage{xcolor}
\usepackage[framemethod=TikZ]{mdframed}
\usepackage{amsthm}
\usepackage{thmtools}
\usepackage{thm-restate}
\usepackage{hyperref}
\usepackage{cleveref}
% NEW: etoolbox is needed to patch the proof environment
\usepackage{etoolbox}

% Definir colores personalizados para las cajas
\definecolor{mygrayback}{RGB}{245, 245, 245}
\definecolor{mygrayframe}{RGB}{80, 80, 80}
\definecolor{mygrayframeproof}{RGB}{140, 140, 140}

% Simbolo demostración
\renewcommand{\qedsymbol}{$\blacksquare$}

% Estilo de caja principal para teoremas y otros entornos
\declaretheoremstyle[
  headfont=\bfseries,
  bodyfont=\normalfont,
  mdframed={
    backgroundcolor=mygrayback,
    linecolor=mygrayframe,
    innertopmargin=8pt,
    innerbottommargin=8pt,
    skipabove=1em,
    skipbelow=1em,
    roundcorner=5pt
  }
]{thmbox}

% Definición de entornos (todos usan el estilo thmbox)
\declaretheorem[numberwithin=section, style=thmbox, name=Teorema]{theorem}
\declaretheorem[sibling=theorem, style=thmbox, name=Lemma]{lemma}
\declaretheorem[sibling=theorem, style=thmbox, name=Proposição]{proposition}
\declaretheorem[sibling=theorem, style=thmbox, name=Observação]{note}
\declaretheorem[sibling=theorem, style=thmbox, name=Corolário]{corollary}
\declaretheorem[sibling=theorem, style=thmbox, name=Exemplo]{example}
\declaretheorem[sibling=theorem, style=thmbox, name=Definição]{definition}
\declaretheorem[sibling=theorem, style=thmbox, name=Notação]{notation}

% Configuración de mdframed para el entorno proof
\renewenvironment{proof}[1][\proofname]{%
  \begin{mdframed}[
    frametitle={#1},
    frametitlerule=true,
    frametitlerulecolor=black,
    frametitlerulewidth=0.2pt,
    linecolor=black,
    backgroundcolor=white,
    roundcorner=5pt,
    innertopmargin=6pt,
    innerbottommargin=12pt,  % más espacio abajo
    innerrightmargin=12pt,   % más espacio derecha
    skipabove=1em,
    skipbelow=1em
  ]%
  \pushQED{\qed} % <-- activa el símbolo
}{%
  \popQED          % <-- imprime el símbolo dentro del marco
  \end{mdframed}
}

% Configuración cleveref para los nuevos entornos
\crefname{theorem}{Theorem}{Theorems}
\crefname{lemma}{Lemma}{Lemmas}
\crefname{proposition}{Proposition}{Propositions}
\crefname{note}{Note}{Notes}
\crefname{corollary}{Corollary}{Corollaries}
\crefname{example}{Example}{Examples}
\crefname{definition}{Definition}{Definitions}



